\section{Architekturmuster}
Gewählt wurde das Server Client Archtuekturmster. Dieses beschreibt das Zusammenarbeiten eines Servers, der Dienste bereitstellt, und eines oder mehrerer Clients, die diese Dienste nutzen. In unserem Fall ist der Server für die Verwaltung der Turnierdaten und die Bereitstellung von APIs verantwortlich, während die Clients (z.B. Webbrowser) die Benutzeroberfläche darstellen und mit dem Server kommunizieren, um Daten zu senden und zu empfangen. Dieses Muster ermöglicht eine klare Trennung von Verantwortlichkeiten, erleichtert die Wartbarkeit und Skalierbarkeit der Anwendung.
\subsection{Das Server-Client-Modell}

Das Server-Client-Modell ist ein grundlegendes Architekturmuster, das die Interaktion zwischen einem zentralen Server und mehreren Clients beschreibt. Der Server stellt Ressourcen oder Dienste bereit, während die Clients diese Ressourcen anfordern und nutzen. In unserem Projekt fungiert der Server als zentrale Datenbank und Logikschicht, die die Turnierdaten verwaltet und APIs bereitstellt, um diese Daten zu manipulieren. Die Clients, in Form von Webbrowsern oder mobilen Anwendungen, greifen auf diese APIs zu, um Informationen anzuzeigen und Benutzereingaben zu verarbeiten. Dieses Modell ermöglicht eine effiziente Verteilung von Aufgaben und eine klare Trennung von Verantwortlichkeiten zwischen Server und Client.

\subsection{Unser Server Client-Modell}

In diesem Fall wurde zumeist ein reines Server-Client-Modell verwendet, da dies die Anforderungen der Anwendung Vollständig erfüllt. Die einzige Abänderung, welche vorgenommen wurde ist die Kommunikation, welche nicht wie gewohnt über das Internet erfolgt, sondern über eine lokale Schnittstelle, da die Anwendung für den Einsatz auf einem lokalen Gerät konzipiert ist. Diese Anpassung ermöglicht eine schnellere und stabilere Kommunikation zwischen Server und Client, da sie nicht von externen Netzwerkbedingungen abhängig ist. Trotz dieser Änderung bleibt die grundlegende Struktur des Server-Client-Modells erhalten, da der Server weiterhin die zentrale Rolle bei der Verwaltung der Daten und der Bereitstellung von Diensten spielt, während die Clients die Benutzeroberfläche darstellen und mit dem Server interagieren.

